\documentclass[10pt,journal,compsoc]{IEEEtran}

\usepackage[utf8]{inputenc}
\usepackage[T1]{fontenc}
\usepackage{textcomp}
\ifCLASSOPTIONcompsoc
  \usepackage[nocompress]{cite}
\else
  \usepackage{cite}
\fi
\usepackage{amsmath,amssymb}
\usepackage{booktabs}
\usepackage{array}
\usepackage{graphicx}
\usepackage{xcolor}
\usepackage{pifont}
\usepackage{url}
\usepackage[hidelinks]{hyperref}
\def\UrlBreaks{\do\/\do-\do\_\do\.\do\:}
\Urlmuskip=0mu plus 1mu\relax
\renewcommand{\arraystretch}{1.1}

\title{GOV-AR: Risk-Bounded Multi-Tenant LLM Admission and Routing\\
under Delayed Token-Cost Feedback}

\author{Anonymous Authors}

\begin{document}

\IEEEtitleabstractindextext{%
\begin{abstract}
Enterprise LLM gateways increasingly share model endpoints across multiple
tenants, yet the final output-token cost of an individual request is only known
after completion and settlement may arrive late. This creates an online control
problem that is not addressed by post-hoc FinOps reporting alone. We present
GOV-AR, a reproducible implementation path for governance-aware admission, routing, and reservation
under delayed token-cost feedback. GOV-AR combines a hard governance filter, a
reservation policy over predicted output lengths, tenant-isolated settled versus
reserved budget accounting, and first-class \texttt{admit}, \texttt{queue},
\texttt{reject}, and \texttt{abstain} actions. The method is implemented as a
research workspace plus an operator-side admission service aligned with an
existing Kubernetes operator for AI FinOps and sovereignty. We verify the
implementation with trace-driven experiments
covering delayed settlement, multi-tenant contention, drift response, fault
injection, scalability, live Azure validation, and component ablation. In the
current evidence, quantile reservation achieves higher throughput than
mean-plus-standard-deviation reservation in one delayed-settlement trace, while
mean-plus-standard-deviation avoids overshoot in other cases at the cost of
higher slack. We also validate a local Docker/Kind/Helm execution path and live
Azure OpenAI and Foundry invocations using synthetic prompts. The result is a
reproducible experimental basis for a fuller operator-integrated GOV-AR system,
including an optional PostgreSQL-backed ledger and a proxy-triggered
admit/settle/cancel path, with explicit documentation of remaining gaps for
literature hardening, deeper Envoy-native integration, and final publication
packaging.
\end{abstract}

\begin{IEEEkeywords}
LLM gateways, admission control, delayed feedback, Kubernetes operators,
multi-tenant systems, FinOps, sovereign AI, Azure OpenAI.
\end{IEEEkeywords}}

\maketitle
\IEEEdisplaynontitleabstractindextext
\IEEEpeerreviewmaketitle

\section{Introduction}
Large-language-model gateways now act as shared infrastructure for multiple
applications, teams, and tenants. In regulated enterprises, each routing
decision must respect cost ceilings, provider and residency constraints, and
application-specific guardrails. However, the output token count of a request is
unknown at admission time, multiple requests can be in flight concurrently, and
cost telemetry may arrive only after delayed settlement. These properties make
budget safety an online systems problem rather than a purely retrospective
reporting problem.

The operator in this repository already provides telemetry ingestion, cost
attribution, budget phase tracking, sovereignty enforcement, and route
recommendation. What it lacks is a first-class reserve-dispatch-settle method
for uncertain, concurrent, and delayed-cost requests. GOV-AR addresses that
gap.

\paragraph{Contributions.}
This manuscript makes three scoped contributions consistent with the current
state of implementation:
\begin{enumerate}
  \item a formulation of multi-tenant LLM admission and routing under delayed
        token-cost feedback with settled and reserved budget views;
  \item a reproducible GOV-AR scaffold with governance filtering, reservation,
        idempotent settlement, queueing, abstention, and delayed-settlement
        replay;
  \item an evidence package spanning E0--E7 experiments, local Docker/Kind/Helm
        execution, and live Azure OpenAI plus Foundry validation with synthetic
        prompts.
\end{enumerate}

\section{Related Work}
Adaptive LLM routing under budget constraints has been studied in recent work
such as Adaptive LLM Routing under Budget Constraints~\cite{panda2025adaptive}.
Cost-aware cascades and routers such as FrugalGPT~\cite{chen2023frugalgpt},
RouteLLM~\cite{ong2024routellm}, HYBRID LLM~\cite{ding2024hybrid}, and
RouteLMT~\cite{luo2026routelmt} show
that query-dependent model selection can materially improve cost-quality
trade-offs. Recent conformal routing work~\cite{uddin2026conformal} adds
distribution-free safety guarantees to the routing decision itself.
Recent ACL 2026 work further expands the routing design space: InferenceDynamics
targets adaptable routing across broad model pools, FLARE emphasizes
length-aware resource efficiency, and MTRouter extends cost-aware routing to
multi-turn trajectories~\cite{shi2026inferencedynamics,fu2026flare,zhang2026mtrouter}.
Delayed feedback and online resource-constrained decision-making also have
deeper theoretical roots in online learning under delayed feedback
~\cite{joulani2013delayed} and contextual bandits with knapsack-style resource
constraints~\cite{badanidiyuru2014resourceful,han2023cbwk}. Output-length-aware
serving further motivates prediction-based reservation and unknown-length-aware
scheduling~\cite{qiu2024ssjf,xie2026forelen,ruan2026libra}. The gap targeted here is their combination
with tenant-isolated budget liability, explicit queueing and abstention
decisions, hard governance filtering, and a Kubernetes-native control-plane
execution context. None of the verified primary sources above combines delayed
settlement, in-flight financial liability, tenant-isolated reservation
accounting, and operator-integrated governance enforcement in one end-to-end
method.

\section{Problem Formulation}
For each incoming request $r$ from tenant $t$, GOV-AR chooses one of
\texttt{admit(model,reservation)}, \texttt{queue}, \texttt{reject}, or
\texttt{abstain}. Hard governance constraints filter infeasible providers and
models before soft scoring begins. At decision time, the system knows the
tenant, request metadata, historical output-length samples, and current
settled/reserved budget state, but not the final output token count nor the
ultimate settled cost.

The key accounting abstraction is:
\[
\mathrm{available}_t(k)=\mathrm{budget}_t-\mathrm{settled}_t(k)-\mathrm{reserved}_t(k).
\]

This separates already committed spend from in-flight liability. A request can
only be dispatched after its reservation is atomically committed.

\section{GOV-AR Design}
The current GOV-AR scaffold has four core pieces.

\paragraph{Governance filtering.}
Candidates that violate hard policy are removed before budget or utility are
considered.

\paragraph{Reservation policy.}
The scaffold currently supports empirical quantile reservation and mean plus
standard deviation reservation over predicted output token samples.

\paragraph{Tenant ledger.}
Each tenant tracks both settled and reserved cost. A per-request ledger stores
reservation amount, model, status, and settlement event identifier. Settlement
is idempotent for repeated delivery of the same event and rejects conflicting
duplicate events.

\paragraph{Replay engine.}
The replay harness supports delayed settlement, queue retries, drift-aware
abstention, and request-level event traces.

\section{Implementation}
The implementation is intentionally split between an operator audit of the
existing platform and a research workspace in \texttt{article3/}. The workspace
contains:
\begin{itemize}
  \item \texttt{src/admission}: the decision layer with \texttt{admit},
        \texttt{queue}, \texttt{reject}, and \texttt{abstain};
  \item \texttt{src/ledger}: tenant budget state plus per-request reservation
        records and idempotent settlement;
  \item \texttt{src/predictor}: reservation bounds and simple drift detection;
  \item \texttt{src/trace\_gateway}: trace generation, replay, and metrics;
  \item \texttt{cmd/experiment}: runners for E0--E7;
  \item \texttt{infra/}: local Docker, Kind, Helm, and Azure live-validation
        automation.
\end{itemize}

On the operator side, the implementation now also includes an optional
\texttt{gov-ar-admission} service that reads \texttt{AIBudgetPolicy},
\texttt{AIRoutingPolicy}, \texttt{AIModel}, and \texttt{AIProvider}, exposes
\texttt{/v1/admit}, \texttt{/v1/settle}, \texttt{/v1/cancel}, and
\texttt{/v1/liability/\{tenant\}}, and can persist reservations and settlements
either in memory or in PostgreSQL. A proxy-triggered live path in the existing
\texttt{header-proxy} can call \texttt{admit} before upstream dispatch and
\texttt{settle} or \texttt{cancel} after response completion. The local
infrastructure now includes dedicated experiment and admission-service
container images, a Helm chart that runs the experiment binary as a Kubernetes
Job, and Kind automation that successfully deployed and executed the job in a
dedicated local cluster.

\section{Evaluation}
The present evaluation is a reproducible evidence package rather than a final
publication-grade study. Still, every numeric claim in this section is tied to
a generated artifact under \texttt{article3/experiments/}.

\subsection{E1 and E2: Delayed Settlement and Multi-Tenancy}
In the deterministic E1 delayed-settlement trace, quantile reservation admitted
two requests and settled a total cost of 12.0 with 2.0 slack and no overshoot.
The mean-plus-standard-deviation variant admitted one request, queued one, and
settled 7.0 with 1.0 slack and no overshoot. In the deterministic E2
multi-tenant trace, both variants admitted three requests and queued one, but
the quantile variant incurred 0.5 overshoot while mean-plus-standard-deviation
used more slack (1.5) and avoided overshoot.

\subsection{E3 and E4: Drift and Faults}
E3 demonstrates conservative fallback under drift. With a mean ratio drift of
2.08 over the configured threshold, the scaffold admitted a non-drifted
fallback candidate for one request and abstained on a second request when no
safe candidate remained. E4 validates that the ledger accepts the first
settlement event, tolerates replay of the same event idempotently, rejects a
conflicting second event, releases reservation on expiry, and abstains when the
decision layer has insufficient fresh evidence.

\subsection{E5: Scalability}
The E5 scalability sweep replays synthetic traces from 50 to 1000 requests. The
largest current run completes in sub-millisecond to millisecond-scale local
replay time, with throughput in the hundreds of thousands of requests per
second in this lightweight harness. These numbers characterize the replay
engine, not an end-to-end production gateway.

\subsection{E6: Azure Live Validation}
E6 validates live infrastructure and real deployments already present in the
subscription. The experiment discovered four cognitive-service resources and
three relevant deployments: \texttt{gpt-france-mini} in
\texttt{francecentral}, \texttt{gpt-us-mini} in \texttt{eastus}, and
\texttt{mistral-large-latest} in a Foundry/AIServices account in
\texttt{westus2}. Live calls with a synthetic prompt returned HTTP 200 and the
expected text \texttt{OK} on both Azure OpenAI deployments and on the Foundry
Mistral deployment, with token-usage and latency metadata captured in the raw
artifact.

\subsection{E7: Ablation}
E7 compares quantile, mean-plus-standard-deviation, an aggressive low
reservation baseline, and a no-queue policy. In the current budget-delay
scenario, the aggressive low reservation variant yields higher settled total
than conservative reservation but incurs overshoot, while disabling queueing
turns deferred decisions into outright rejects. This supports the claim that
both reservation policy and queue semantics materially shape the
safety-throughput trade-off.

\section{Discussion}
The strongest result of the current artifact is reproducibility across several
axes at once: deterministic trace replay, live Azure validation, local
containerization, cluster execution, and an operator-side admission service
with a persistent ledger option. The weakest aspect is still depth of
integration into the final operator control path. GOV-AR now includes a
proxy-triggered reserve-settle path, but it is not yet an Envoy-native or
controller-driven production integration. The literature review and manuscript
also remain less mature than the implementation and experiment harness.

\section{Threats to Validity}
The main threats are external validity and maturity bias. Synthetic traces do
not capture the full burstiness or semantic variability of production prompts.
Replay throughput does not equal production gateway throughput. Azure live
validation used only short synthetic prompts and existing deployments; it does
not by itself prove deployment-scale economics. Finally, the related-work
corpus is materially stronger than in the initial draft, but the broader review
still needs a final saturation pass before submission freeze.

\section{Conclusion}
GOV-AR frames enterprise LLM routing as a joint admission, reservation, and
settlement problem under delayed token-cost feedback and hard governance
constraints. The current artifact already demonstrates a reproducible decision
kernel, a local Docker/Kind/Helm execution path, an operator-side admission
service with optional PostgreSQL persistence, a proxy-triggered live
admit/settle path, and successful Azure live validation against real Azure
OpenAI and Foundry deployments. The next required steps are deeper Envoy-native
operator integration, a final saturation pass on the literature review, and
completion of the final paper and replication package.

\bibliographystyle{IEEEtran}
\bibliography{references}
\end{document}
